\chapter{The Equivalence Principle}

\section{Gravitational mass and inertial mass.}

\begin{points}
\item In Newtonian mechanics mass plays a fundamental role, since, through Newton's second law, it determines the acceleration that a body undergoes when it is subjected to a force,
\[
  \vec F=m\vec a .
\]
This parameter is called the inertial mass of the body.

\item On the other hand, the law of universal gravitation states that the gravitational force between two bodies is proportional to their mass.

\item From a conceptual point of view, the mass appearing in the law of universal gravitation is completely different from the previous one. In Newton's law it is a ``resistance'' of the particle to changing its velocity, whereas in gravitation it corresponds to the value of the gravitational ``charge'' of the particle, i.e.\ the property of the object that generates, and is affected by, the gravitational interaction.

\item Given the conceptual difference between the two, it is useful to distinguish them and to call \emph{gravitational mass} the mass term appearing in the law of universal gravitation.

\item From this point of view, we have two physical quantities that are a priori completely different and that, from an observational point of view, turn out to have the property of being proportional to each other, with a universal proportionality factor (which can therefore be chosen equal to one by a suitable choice of units).
\end{points}

\section{Galileo's equivalence principle.}

\begin{points}
\item This experimental fact is of such a fundamental nature that it can be raised to a principle of physics. This principle is known as ``Galileo's equivalence principle'' or also the ``weak equivalence principle''.

\item As with every physical principle, we must ask ourselves to what extent it has been tested experimentally.

\item Galileo Galilei (1564--1642) discovered that bodies fall at a rate independent of their mass. The usual way of describing these experiments is through the famous (and very probably apocryphal) experiment of Galileo dropping objects of different composition and density from the top of the Leaning Tower of Pisa to check that they all take the same amount of time to reach the ground.

\item The experiment that Galileo did perform consisted in letting two heavy objects roll down an inclined plane with a small angle of inclination, so that the fall time was long and the motion of objects of different compositions and sizes could thus be compared.

\item This experiment, together with improved versions by Huygens (1629--1695), allowed Newton to conclude that the force of gravity is proportional to the (inertial) mass of the object that falls. By Newton's third law one can conclude that the gravitational force is also proportional to the (inertial) mass of the object that creates the force.

\item However, Newton was fully aware that this hypothesis could be only approximate and that, therefore, the inertial mass of the second law could be different from the gravitational mass.

\item Distinguishing between the two by means of the symbols $m_i$ and $m_g$ respectively, and calling $\vec g$ the gravitational field generated by any (gravitational) mass distribution, the acceleration of a point particle with masses $m_i$ and $m_g$ is given by
\[
  \vec a=\frac{m_g}{m_i}\vec g .
\]
The experimental hypothesis to be tested is whether the ratio $m_g/m_i$ can be different for different bodies.

\item Under these conditions, a pendulum of length $\ell$ performing small oscillations has a period
\[
  T=2\pi\sqrt{\frac{m_i\ell}{m_gg}} .
\]

\item By taking pendulums of equal length and different compositions, one can compare the periods and see to what extent they coincide.

\item Newton carried out this kind of experiment and observed no differences in the periods of the pendulums.

\item These experiments were improved by Bessel in 1830, also with negative results for the difference in the periods.

\item However, the measurements of the periods of pendulums are not precise enough (essentially because of the damping of the pendulum), so the equality between inertial and gravitational mass could be established with a precision of at most $10^{-3}$.

\item In 1889 Loránd von Eötvös, of Budapest, devised a much more precise method to test experimentally the equality between inertial and gravitational mass. His method made it possible to conclude that the ratio $\frac{m_i}{m_g}$ for two different substances does not differ by more than one part in $10^9$, which is a really very high precision.

\item The method consists in using the rotation of the Earth and the centrifugal acceleration that it generates in a reference frame attached to the laboratory.

\item The experimental set-up is the following: two equal masses joined by a rod 40\,cm long are suspended.

\item The following external forces act on the rod:

(i) The force of gravity $\vec F_g=m_g\vec g$ on each of the particles $A$ and $B$.

(ii) The centrifugal force on each of the particles. Being an inertial force, it is proportional to the inertial mass of the particles,
\[
  \vec F_C=-m_i\vec\omega\times(\vec\omega\times\vec r),
\]
where $\vec\omega$ is the angular velocity of rotation of the Earth and $\vec r$ is the position vector of the particle with respect to the centre of the Earth. Writing $\vec r=\vec R_C+\vec r\,'$, where $\vec R_C$ is the position vector of the suspension point $C$ with respect to the centre of the Earth and $\vec r\,'$ the position of the particle relative to the point $C$, we can approximate the centrifugal force by
\[
  \vec F_C=-m_i\vec\omega\times\left(\vec\omega\times\vec R_C\right).
\]

(iii) The force exerted by the suspension fibre of the rod, which is directed along the fibre itself.

(iv) In addition, because of the torsion of the fibre, the latter exerts a torque on the rod, directed along the fibre itself.

\item We place ourselves at the equilibrium position of the system, in which the rod is horizontal. We call $\vec e_1$ the unit vector at $C$ along the fibre, and $\vec e_2$ the unit vector at $C$ pointing from particle $A$ towards $B$.

\item We define $\vec g_C:=\vec\omega\times\left(\vec\omega\times\vec R_C\right)$, so that the centrifugal force on each particle is $\vec F_C=-m_i\vec g_C$.

\item We call $m_i^A$ and $m_i^B$ the inertial masses of the particles $A$ and $B$, and $m_g^A$, $m_g^B$ their gravitational masses.

\item The equilibrium condition for the forces and for the torques imposes the conditions
\begin{gather*}
  T\vec e_1+(m_g^A+m_g^B)\vec g-(m_i^A+m_i^B)\vec g_C=0,\\
  M\vec e_1+l_B\vec e_2\times\left(m_g^B\vec g-m_i^B\vec g_C\right)-l_A\vec e_2\times\left(m_g^A\vec g-m_i^A\vec g_C\right)=0
\end{gather*}
where $l_A$ is the length between the point $C$ and the point $A$, and similarly for $l_B$.

\item We solve the first equation to obtain $\vec g$:
\[
  \vec g=\frac{-T}{m_g^A+m_g^B}\vec e_1+\frac{m_i^A+m_i^B}{m_g^A+m_g^B}\vec g_C
\]
and substitute it into the torque-balance equation. Using
\begin{align*}
  m_g^B\vec g-m_i^B\vec g_C&=\frac{-Tm_g^B}{m_g^A+m_g^B}\vec e_1+\frac{m_g^Am_g^B}{m_g^A+m_g^B}\left(\frac{m_i^A}{m_g^A}-\frac{m_i^B}{m_g^B}\right)\vec g_C,\\
  m_g^A\vec g-m_i^A\vec g_C&=\frac{-Tm_g^A}{m_g^A+m_g^B}\vec e_1-\frac{m_g^Am_g^B}{m_g^A+m_g^B}\left(\frac{m_i^A}{m_g^A}-\frac{m_i^B}{m_g^B}\right)\vec g_C,
\end{align*}
we obtain
\[
  M\vec e_1+\frac{T}{m_g^A+m_g^B}\left(l_Bm_g^B-l_Am_g^A\right)\vec e_3+(l_A+l_B)\frac{m_g^Am_g^B}{m_g^A+m_g^B}\left(\frac{m_i^A}{m_g^A}-\frac{m_i^B}{m_g^B}\right)\vec e_2\times\vec g_C=0
\]

\item We conclude that the torque exerted by the fibre must be proportional to the quantity
\[
  \left(\frac{m_i^A}{m_g^A}-\frac{m_i^B}{m_g^B}\right).
\]
More precisely,
\[
  M=\mu_g\left(\frac{m_i^A}{m_g^A}-\frac{m_i^B}{m_g^B}\right)\vec e_1\cdot\left(\vec g_C\times\vec L\right)
\]
where $\vec L$ is the position vector of the mass $B$ with respect to $A$, and $\mu_g$ is the (gravitational) reduced mass of the system.

\item By measuring the deflection angle of the rod for different masses, we can determine $\left(\frac{m_i^A}{m_g^A}-\frac{m_i^B}{m_g^B}\right)$.

\item Eötvös's original experiment of 1889 was improved by Eötvös, Pekár and Fekete in 1922. The precision achieved was
\[
  \left|\frac{m_i-m_g}{m_g}\right|\leq3\cdot10^{-9}.
\]
The experiment was repeated by Dicke, Krotkov and Roll in 1964, using the orbital motion of the Earth around the Sun, instead of the rotation of the Earth, as the source of the field of inertial forces. They obtained
\[
  \left|\frac{m_i-m_g}{m_g}\right|\leq1.5\cdot10^{-11}.
\]
More recent measurements have provided the following bounds:
\begin{align*}
  &\text{Adelberger \emph{et al} (1990)}, && \frac{|\Delta m|}{m}=\left\{\begin{array}{ll}(0.2\pm1.0)\cdot10^{-11}&\quad\text{Cu / Be}\\(0.5\pm1.3)\cdot10^{-11}&\quad\text{Al / Be}\end{array}\right.\\
  &\text{Baesler \emph{et al} (1999)}, && \frac{|\Delta m|}{m}\leq5\cdot10^{-13}
\end{align*}

There are several space missions under way or being designed to test the weak equivalence principle with higher precision. In April 2016 the MicroSCOPE satellite was launched with the aim of establishing a bound of $10^{-15}$. In December 2017 the first results of this mission were published, which indeed confirm the bound of $10^{-15}$ on the validity of the weak equivalence principle. The launch of a joint NASA/ESA satellite called STEP (Satellite Test of the Equivalence Principle) is also planned (with no specific date), with the mission of measuring $\Delta m/m$ with a precision of $10^{-18}$.
\end{points}

\section{Gravitational forces and inertial forces.}

\begin{points}
\item Galileo's equivalence principle thus has a very good experimental basis.

\item What consequences can we draw from this principle?

\item If we consider a Cartesian reference frame $S$ that moves (without rotating) with acceleration $\vec a_0(t)$ with respect to an inertial one, and we want to write Newton's laws in the frame $S'$, we must include a non-inertial force $\vec F_I=-m_i\vec a$, so that Newton's law in $S$ is
\[
  m_i\vec a=\vec F-m_i\vec a_0
\]
where $\vec F$ is the set of applied forces. On the other hand, the motion of a particle with respect to an inertial reference frame, on which a uniform gravitational force acts in addition to the applied forces $\vec F$, is given by Newton's law
\[
  m_i\vec a=\vec F+m_g\vec g .
\]
The two situations are completely indistinguishable from the dynamical point of view provided that $\vec g=-\vec a_0$.

\item We conclude that in Newtonian mechanics the laws of dynamics are equivalent for a reference frame with constant acceleration (with respect to an inertial one) and for an inertial reference frame with a constant and uniform gravitational field.

\item Let us recall Galileo's principle of relativity: \emph{It is impossible to determine by mechanical processes whether a reference frame is at rest or in uniform rectilinear motion}. In other words, the laws of mechanics are equivalent in a given inertial reference frame and in any other reference frame moving with constant velocity with respect to the former.

\item The weak equivalence principle seems to indicate that accelerated reference frames and inertial reference frames with gravitational fields could be indistinguishable by means of mechanical processes.

\item However, the previous equivalence relied heavily on the fact that $\vec g$ is uniform.

\item If we consider an arbitrary gravitational field $\vec g(\vec x,t)$, it will not be possible to put a Cartesian reference frame $S$ in a state of motion such that the inertial force corresponds to $m\vec g$.

\item However, we can consider a reference frame $S'$ that does not rotate with respect to the previous one and whose origin $O'$ follows a free-fall trajectory in the field $\vec g(\vec x,t)$. If we consider a particle on which no force acts other than gravity, and which is located near $O'$ and initially moving with velocity $\vec v$ with respect to $O'$, we can ask ourselves what motion $S'$ will see.

\item The answer is that, as long as the particle does not move too far from the origin $O'$ (compared with the scales of inhomogeneity of $\vec g(\vec x,t)$), $O'$ will see the particle moving with constant velocity $\vec v$.

\item For example, if in a freely falling satellite we let go of a pen, it will move away from us with practically constant velocity. Although the gravitational field in which it is located is not uniform, on the scale of the satellite the variations of the field are very small and the effects of these inhomogeneities are negligible.

\item What we have achieved is to cancel the effect of the gravitational field by letting ourselves fall in it. If we had no information about the kind of frame we were in, we could argue that, since we observe that free particles do not accelerate, our reference frame is inertial and no external gravitational field acts.

\item Naturally, this cancellation of the gravitational field is only local, both in space and in time.

\item For example, let us consider an elevator in free fall in the gravitational field of the Earth. If we are inside it we will not notice any weight. We can ask ourselves what happens to two particles that we release from rest at the same height (inside the elevator). Since both particles move attracted by the Earth, their motion will be radial towards the centre of the Earth. Seen from the elevator, we will see that the two particles approach each other very slowly (and this attraction has nothing to do with their own mutual gravitational attraction, which we are neglecting completely). This is a tidal force, which corresponds to variations of the gravitational field.

\begin{carroll}[: why only locally]
In a sealed box accelerated by a rocket in empty space, free particles released from rest always ``fall'' along parallel straight lines, with the same acceleration. In a large box at rest (or freely falling) in the gravitational field of the Earth, they move instead towards the centre of the Earth, i.e.\ along \emph{converging} directions. By observing such relative motions an observer inside a sufficiently large box could distinguish a gravitational field from a uniform acceleration. This is why the equivalence between the two only holds in sufficiently small regions of space-time, where these inhomogeneities (tidal effects) cannot be detected.
\begin{center}
\begin{tikzpicture}[scale=0.85,>=Stealth]
  \begin{scope}
    \draw[thick] (0,0) rectangle (4,2.2);
    \foreach \x in {1,3}{\fill (\x,1.7) circle (3pt); \draw[->,thick] (\x,1.5) -- (\x,0.4);}
    \foreach \x in {1,3}{\draw[thick] (\x-0.35,0) arc (180:360:0.35 and 0.3);
      \draw[orange!80!red,thick,decorate,decoration={zigzag,amplitude=2pt,segment length=4pt}] (\x-0.25,-0.35) -- (\x,-0.9) -- (\x+0.25,-0.35);}
    \node[below,align=center] at (2,-1.1) {accelerated box (rocket):\\ parallel motion};
  \end{scope}
  \begin{scope}[shift={(6.5,0)}]
    \draw[thick] (0,0.6) rectangle (4,2.8);
    \fill (0.7,2.3) circle (3pt); \fill (3.3,2.3) circle (3pt);
    \draw[->,thick] (0.75,2.1) -- (1.15,1.0);
    \draw[->,thick] (3.25,2.1) -- (2.85,1.0);
    \draw[thick] (-0.3,-1.3) arc (150:30:2.66);
    \node at (2,-0.55) {Earth};
    \node[below,align=center] at (2,-1.3) {box in a gravitational field:\\ converging motion};
  \end{scope}
\end{tikzpicture}
\end{center}
\end{carroll}

\item We conclude that the reference frame is inertial (for mechanical processes) on scales on which tidal forces are negligible.
\end{points}

\section{Einstein's equivalence principle.}

\begin{points}
\item One of Einstein's great merits was to raise these mechanical properties, which are a consequence of Newton's laws and of Galileo's equivalence principle, to a principle of physics (and therefore conceptually prior to Newton's laws), and furthermore to extend it to all the laws of physics. This is a very important conceptual leap, which was not based on any previous experiment.
\end{points}

\textbf{Principle 1} (Einstein's equivalence principle)\textbf{.} \emph{All the laws of physics (except gravitation) are the same, on sufficiently local scales, in a freely falling reference frame in the presence of a gravitational field or in an inertial reference frame in the absence of gravity.}

An alternative version of this principle is:

\emph{The laws of physics (except gravitation) are the same, on sufficiently local scales, in an inertial reference frame with a gravitational field and in an accelerated reference frame.}

\begin{points}
\item This principle turned out to be the key for Einstein to be able to develop a theory of gravitation incorporating the essential principles of special relativity.

This process was not easy. Einstein published the equivalence principle in 1907 and it took him 8 years to develop his theory of gravity, which was named the theory of general relativity.

\item Part of the difficulty (but not all of it) was that Einstein had to learn the appropriate mathematics in which the new theory had to be formulated. In this task he had the help of a colleague and friend, the mathematician Marcel Grossmann.
\end{points}

\section{Space-time geometry and gravity.}

\begin{points}
\item Before entering the theory proper, we can draw some basic consequences of Einstein's equivalence principle.

\item Let us consider a ray of light moving in an inertial reference frame without gravitational fields. Its trajectory will be a straight line travelled at constant speed $c$.

\item Let us consider first a merry-go-round rotating with constant angular velocity about an axis parallel to the ray of light. This frame is non-inertial and the trajectory it observes for the ray of light is a helix. This is a curved trajectory. Since the ray moves very fast, under normal conditions we will not be able to observe this curvature. To see a substantial piece of the helix (and therefore notice the curvature of the trajectory), the frame must rotate at a sufficiently high speed. More precisely, $\omega D$ must be comparable with $c$, where $D$ is the distance between the trajectory of the light and the axis of the merry-go-round.

\item By the (Einstein) equivalence principle we conclude that a ray of light in the presence of a gravitational field follows a curved trajectory.

\item We can describe another example of the same phenomenon in the following way. Let us consider an elevator with a small slit on one side through which a ray of light enters. The elevator is in free fall in the gravitational field of the Earth. By the equivalence principle, the reference frame of the elevator is equivalent to an inertial one, so the ray of light must move in a straight line with respect to the elevator. For this to happen, the trajectory of the ray as seen from the Earth must bend downwards. We reach the same conclusion as before: the gravitational field bends the trajectories of light rays.

\item Let us consider again an inertial reference frame $S$ and another one $S'$ that rotates with respect to the former, and let us assume that there is no gravitational field (with respect to the frame $S$). We choose the $z$ and $z'$ axes along the rotation axis.

\item Let us take a circle centred at the origin (formed, for example, by a thin tube) and in the plane perpendicular to the $z$ axis. The inertial reference frame measures a length $\ell$ of the tube and a diameter $d$ of the circle, such that the ratio $\frac\ell d$ has the universal value
\[
  \frac\ell d=\pi .
\]

\item We can ask ourselves what the observer in the reference frame $S'$ measures. By symmetry, this observer must still see a circle (all its points must be at the same distance from the origin $O'$). To measure lengths, it must use rulers at rest in its own reference frame. The rulers used to measure the diameter of the circle move (with respect to the inertial frame $S$) in a direction perpendicular to the ruler, so that according to special relativity they have the same length both in $S$ and in $S'$. We conclude that $S'$ still measures $d$ as the diameter of the circle.

\item On the other hand, the rulers used to measure the perimeter of the circle move with longitudinal velocity $\frac{\omega d}{2}$, so that according to $S$ they have a length
\[
  L=L_0\sqrt{1-\frac{v^2}{c^2}}=L_0\sqrt{1-\frac{d^2\omega^2}{4c^2}},\qquad L_0=\text{proper length of the ruler}
\]
because of the length contraction of special relativity. Therefore, for them to cover the whole length of the circle, $S$ asserts that the number of rulers that must be placed is $n=\frac{\ell}{L_0\sqrt{1-\frac{d^2\omega^2}{4c^2}}}$. This is a natural number that counts objects, so $S'$ must count exactly the same number of rulers. Consequently, the perimeter of the circle according to $S'$ is
\[
  \ell'=nL_0=\frac{\ell}{\sqrt{1-\frac{d^2\omega^2}{4c^2}}} .
\]
The ratio between perimeter and diameter is
\[
  \frac{\ell'}{d'}=\frac\ell d\frac{1}{\sqrt{1-\frac{d^2\omega^2}{4c^2}}}=\frac{\pi}{\sqrt{1-\frac{d^2\omega^2}{4c^2}}}>\pi
\]
and we reach the conclusion that the geometry of space according to $S'$ cannot be Euclidean.

\item By the equivalence principle, the presence of gravitational fields also modifies the geometry of space, which can no longer be Euclidean.

\item A third important consequence of the equivalence principle is that the frequency of a photon changes along its trajectory in a gravitational field. This effect is known as the \textbf{gravitational Doppler effect}.

\item We can argue for it by means of the following \emph{Gedankenexperiment}, proposed by Einstein himself.

\item Let us consider a tower of height $h$ on the surface of the Earth, and let us suppose that we drop an object of mass $m$ from the top.

\item When the particle reaches the bottom it has kinetic energy, so that its total energy is greater than its rest energy.

\item Let us suppose that the particle is instantaneously transformed into a photon with the same energy and moving upwards (although this process violates the conservation of linear momentum, we can think that this conversion occurs just as the particle reaches the ground, so that the linear momentum is absorbed by the Earth, in the same way as in the elastic bounce of a ball against a wall).

\item Let us suppose that, when the photon reaches the top of the tower, the inverse process occurs and the photon turns into a material particle at rest (again, the difference in linear momentum can be absorbed by the tower and, ultimately, by the Earth).

\item If the photon did not change its frequency while climbing the tower, the final inertial mass $m'$ would be strictly greater than the initial one, $m'>m$, and we could build a perpetual motion machine of the first kind and violate the conservation of energy, which is one of the pillars of physics.

\item We conclude that the photon must decrease its frequency (energy) when it climbs in the gravitational field.

\item We can make an approximate calculation of how the frequency of the photon must change. The energy of the particle arriving at the bottom is
\[
  \frac{mc^2}{\sqrt{1-\frac{v^2}{c^2}}}=mc^2+\frac12mv^2+O\left(\frac{v^4}{c^4}\right).
\]

\item Using Newton's law, the speed at which it arrives at the bottom is $v=\sqrt{2gh}$, where $h$ is the height of the tower. Therefore $E_f=mc^2+mgh+O(v^4/c^4)$.

\item For energy to be conserved, the photon at the top must have an energy satisfying
\[
  \frac{E'_\gamma}{E^0_\gamma}=\frac{mc^2}{mc^2+mgh+O(4)}=1-\frac{hg}{c^2}+O(4)=\frac{\nu'}{\nu_0}
\]

\item Therefore, in its ascent the photon undergoes a Doppler effect (redshift) described by the following expression
\[
  \nu'=\nu_0\left(1-\frac{gh}{c^2}+O(4)\right).
\]

\item This is a measurable effect. In 1959 Pound and Rebka devised the following experiment (later improved by Pound and Snider in 1965). It consisted in measuring the change in frequency of a photon climbing a column 22.5 metres high (which corresponds to a relative frequency change of $\frac{gh}{c^2}=2.5\cdot10^{-15}$). The experiment was carried out at the Jefferson Laboratory at Harvard. The method used to measure the frequencies consisted in using the Mössbauer effect, i.e.\ the existence of resonances in the absorption and emission of gamma rays by the nuclei of the atoms of a solid. The result of the experiment was the confirmation of the gravitational Doppler effect with a precision of 1\%.

\begin{carroll}[: the gravitational redshift from accelerated boxes, and why space-time is curved]
\textit{(a) Accelerated boxes.} Consider two boxes, a distance $z$ apart, moving far away from any matter with the same constant acceleration $a$. At time $t_0$ the trailing box emits a photon of wavelength $\lambda_0$. The boxes keep a constant distance, so the photon reaches the leading box after $\Delta t=z/c$, during which the boxes have picked up an additional velocity $\Delta v=a\Delta t=az/c$. To first order in $\Delta v/c$, the photon reaching the leading box is therefore redshifted by the ordinary Doppler effect,
\[
  \frac{\Delta\lambda}{\lambda_0}=\frac{\Delta v}{c}=\frac{az}{c^2}.
\]
By the equivalence principle the same must happen in a tower of height $z$ on the surface of a planet, with $a\to g$, which reproduces the result above ($\nu'/\nu_0=1-gh/c^2$). For a non-uniform field one integrates along the path, $\frac{\Delta\lambda}{\lambda_0}=\frac{1}{c^2}\int\partial_z\Phi\,dz=\frac{\Delta\Phi}{c^2}$.
\begin{center}
\begin{tikzpicture}[scale=0.8,>=Stealth]
  % accelerated boxes
  \draw[thick] (0,0) rectangle (0.8,0.6); \draw[thick] (0,2.6) rectangle (0.8,3.2);
  \draw[->] (0.4,3.2) -- (0.4,3.8) node[right] {$a$};
  \draw[->] (0.4,0.6) -- (0.4,1.1) node[right] {$a$};
  \draw[decorate,decoration={snake,amplitude=1.5pt,segment length=5pt},->,red!70!black] (1.1,0.7) -- (1.1,1.8) node[right,black] {$\lambda_0$};
  \draw[<->] (-0.4,0.3) -- (-0.4,2.9) node[midway,left] {$z$};
  \node[below] at (0.4,-0.1) {\small $t=t_0$};
  \begin{scope}[shift={(3.2,1.4)}]
    \draw[thick] (0,0) rectangle (0.8,0.6); \draw[thick] (0,2.6) rectangle (0.8,3.2);
    \draw[decorate,decoration={snake,amplitude=1.5pt,segment length=5pt},->,red!70!black] (1.1,1.6) -- (1.1,2.55);
    \draw[<->] (1.7,0.3) -- (1.7,2.9) node[midway,right] {$z$};
    \node[below] at (0.4,-0.1) {\small $t=t_0+z/c$};
  \end{scope}
  \draw[dashed,gray] (0.8,0.6) -- (3.2,2.0); \draw[dashed,gray] (0.8,3.2) -- (3.2,4.6);
  % tower
  \begin{scope}[shift={(7.5,0)}]
    \fill[pattern=north east lines] (-0.8,-0.3) rectangle (2.4,0);
    \draw (-0.8,0) -- (2.4,0);
    \draw[thick] (0.3,0) -- (0.6,3.0) -- (1.0,3.0) -- (1.3,0);
    \draw (0.38,0.8) -- (1.22,0.8) (0.46,1.6) -- (1.14,1.6) (0.53,2.3) -- (1.07,2.3);
    \draw (0.38,0.8) -- (1.14,1.6) (1.22,0.8) -- (0.46,1.6);
    \draw[thick] (0.35,3.0) rectangle (1.25,3.6);
    \draw[decorate,decoration={snake,amplitude=1.5pt,segment length=5pt},->,red!70!black] (1.9,0.2) -- (1.9,3.4);
    \node[right] at (1.95,0.6) {$\lambda_0$};
    \draw[<->] (-0.4,0) -- (-0.4,3.0) node[midway,left] {$z$};
  \end{scope}
\end{tikzpicture}
\end{center}

\textit{(b) Why space-time must be curved.} A physicist on the ground emits a beam of light from height $z_0$, received at the top of the tower at height $z_1$. The time between the emission of the beginning and of the end of one wavelength is $\Delta t_0=\lambda_0/c$, and the time between the reception of the beginning and of the end is $\Delta t_1=\lambda_1/c$. Since the gravitational field does not change with time, the paths in space-time followed by the leading and trailing edges of the wave must be precisely congruent (a time translation of each other, see the figure). In a flat space-time simple geometry would then tell us that $\Delta t_1=\Delta t_0$. But the gravitational redshift implies $\Delta t_1>\Delta t_0$ (``the clock at the top of the tower appears to run faster''). The fault lies with the ``simple geometry'': a better description is to admit that space-time is curved.
\begin{center}
\begin{tikzpicture}[scale=0.85,>=Stealth]
  \draw[->] (0,0) -- (5,0) node[right] {$z$};
  \draw[->] (0,0) -- (0,4.2) node[above] {$t$};
  \draw (1,0) -- (1,4); \draw (4,0) -- (4,4);
  \node[below] at (1,0) {$z_0$}; \node[below] at (4,0) {$z_1$};
  \draw[thick,->] (1,0.9) .. controls (2,1.1) and (3,1.4) .. (4,1.8);
  \draw[thick,->] (1,2.0) .. controls (2,2.2) and (3,2.5) .. (4,2.9);
  \draw[<->] (0.7,0.9) -- (0.7,2.0) node[midway,left] {$\Delta t_0$};
  \draw[<->] (4.3,1.8) -- (4.3,2.9) node[midway,right] {$\Delta t_1$};
\end{tikzpicture}
\end{center}
\end{carroll}

\item We have already seen that the spatial geometry for an accelerated observer (and therefore also in the presence of a gravitational field) is no longer Euclidean. We also know that, because of special relativity, space and time form a four-dimensional entity called space-time. The space-time of special relativity is described by a metric called the Minkowski metric, which takes the form
\[
  ds^2=-c^2dt^2+dx^2+dy^2+dz^2
\]
where $t,x,y,z$ are the Cartesian time and space coordinates associated with an inertial reference frame.

\item This metric is used to measure space-time intervals between two events. In cylindrical coordinates $\{t,\rho,\phi,z\}$ the Minkowski metric takes the form
\[
  ds^2=-c^2dt^2+d\rho^2+\rho^2d\phi^2+dz^2 .
\]
If we now consider a reference frame rotating with constant velocity around the $z$ axis, a point that initially had coordinates $\{\rho',\phi',z'\}$ and that rotates with constant angular velocity follows the trajectory
\[
  \rho(t)=\rho',\qquad\phi(t)=\phi'+\omega t,\qquad z(t)=z' .
\]
The reference frame $S'$ can decide to coordinatise the events of space-time in such a way that all the points of the previous trajectories are identified by the values $\{\rho',\phi',z'\}$ and by the time $t$ of the clocks in $S$. Therefore the coordinate transformation between the two reference frames is
\[
  \left\{\begin{array}{l}t=t'\\ \rho=\rho'\\ \phi=\phi'+\omega t'\\ z=z'\end{array}\right.
\]

\item If we write the Minkowski metric in the new coordinates we obtain
\begin{align*}
  ds^2&=-c^2dt'^2+d\rho'^2+\rho'^2\left(d\phi'+\omega dt'\right)^2+dz'^2\\
  &=-c^2\left(1-\frac{\rho'^2\omega^2}{c^2}\right)dt'^2+2\omega\rho'^2d\phi'dt'+d\rho'^2+\rho'^2d\phi'^2+dz'^2
\end{align*}
which obviously has a form different from the Minkowski form in the coordinates of the inertial reference frame.

\item By Einstein's equivalence principle we conclude that the space-time geometry of a space-time containing a gravitational field is described by a metric different from the Minkowskian form $-c^2dt^2+dx^2+dy^2+dz^2$.

\item Therefore, we must be able to handle space-times with arbitrary Lorentzian metrics and to derive geometric consequences from them. For this, it is necessary to develop the fundamental tools of differential geometry.
\end{points}
